% در این فایل، دستورها و تنظیمات مورد نیاز آورده شده است.
% بر پایه‌ی commands.tex الگوی رسمی AUTthesis، با افزوده‌های زیر:
%   • بسته‌های listings، tikz، longtable، booktabs و float برای قطعه‌کدها،
%     نمودارهای برداری و جدول‌های بلند.
%   • راه گریز برای قلم‌های PGaramond و IranNastaliq که روی هر سامانه‌ای نصب نیستند.
%-------------------------------------------------------------------------------

% در نسخه‌ی جدید زی‌پرشین، برای تایپ متن‌های ریاضی این سه بسته باید فراخوانی شود.
\usepackage{amsthm,amssymb,amsmath,amsfonts}

% حاشیه‌های صفحه، مطابق دستورالعمل تحصیلات تکمیلی دانشگاه
\usepackage[top=30mm, bottom=30mm, left=25mm, right=30mm]{geometry}

\usepackage{graphicx}
\usepackage{color}
\usepackage{xcolor}
\usepackage{setspace}
\usepackage{titletoc}
\usepackage{tocloft}
\usepackage{enumitem}

% جدول‌ها
\usepackage{array}
\usepackage{longtable}
\usepackage{multirow}
\usepackage{booktabs}

% شکل‌ها
\usepackage{float}

% قطعه‌کدها
\usepackage{listings}

% نمودارهای برداری
\usepackage{tikz}
\usetikzlibrary{positioning,arrows.meta,shapes.geometric,fit,calc,backgrounds}

% لینک‌های رنگی. برای نسخه‌ی نهایی چاپی، خط زیر را غیرفعال و خط پس از آن را فعال کنید.
\usepackage[pagebackref=false,colorlinks,linkcolor=blue,citecolor=red]{hyperref}
%\usepackage[pagebackref=false]{hyperref}

\usepackage[nameinlink]{cleveref}
\AtBeginDocument{%
    \crefname{equation}{برابری}{equations}%
    \crefname{chapter}{فصل}{chapters}%
    \crefname{section}{بخش}{sections}%
    \crefname{appendix}{پیوست}{appendices}%
    \crefname{enumi}{مورد}{items}%
    \crefname{footnote}{زیرنویس}{footnotes}%
    \crefname{figure}{شکل}{figures}%
    \crefname{table}{جدول}{tables}%
    \crefname{theorem}{قضیه}{theorems}%
    \crefname{lemma}{لم}{lemmas}%
    \crefname{corollary}{نتیجه}{corollaries}%
    \crefname{proposition}{گزاره}{propositions}%
    \crefname{definition}{تعریف}{definitions}%
    \crefname{result}{نتیجه}{results}%
    \crefname{example}{مثال}{examples}%
    \crefname{remark}{نکته}{remarks}%
    \crefname{note}{یادداشت}{notes}%
}

\usepackage{fancyhdr}
\usepackage[nottoc]{tocbibind}
\usepackage{makeidx,multicol}
\setlength{\columnsep}{1.5cm}

\usepackage{verbatim}
\makeindex
\usepackage{sectsty}

% فراخوانی زی‌پرشین و تعریف قلم‌ها
\usepackage{xepersian}
\SepMark{-}
\settextfont[Scale=1.2]{B Nazanin}
\setlatintextfont{Times New Roman}
\renewcommand{\labelitemi}{$\bullet$}

%-------------------------------------------------------------------------------
% ارقام فرمول‌ها
%
% الگوی رسمی این خط را فعال دارد تا ارقامِ داخل فرمول‌ها فارسی شود، و برای این کار
% قلم PGaramond لازم است. اگر آن قلم روی سامانه‌ی شما نصب است، خط زیر را فعال کنید.
% در حالت غیرفعال، ارقامِ فرمول‌ها لاتین می‌مانند — که برای متنی با این حجم از
% عبارت‌های عددی و نام‌های لاتین، خواناتر هم هست.
%\setdigitfont[Scale=1.1]{PGaramond}
%-------------------------------------------------------------------------------

% قلم نستعلیق. الگوی رسمی از آن تنها در صفحه‌های «تقدیم» و «سپاس‌گزاری» استفاده
% می‌کند و هر دو صفحه در این نوشتار حذف شده‌اند، پس این تعریف در حال حاضر
% بی‌استفاده است؛ نگه داشته شده تا اگر بعداً یکی از آن صفحه‌ها بازگردد آماده باشد.
\defpersianfont\nastaliq[Scale=2]{IranNastaliq}

\defpersianfont\chapternumber[Scale=3]{B Nazanin}

\renewcommand\proofname{\textbf{برهان}}
\renewcommand{\bibname}{منابع و مراجع}

%-------------------------------------------------------------------------------
% تنظیم قطعه‌کدها
%
% هر قطعه‌کد باید داخل محیط latin قرار گیرد تا جهت آن چپ‌به‌راست بماند. برای
% کوتاه‌شدن کار، محیط‌های آماده‌ی زیر تعریف شده‌اند.
%-------------------------------------------------------------------------------
\definecolor{codebg}{rgb}{0.97,0.97,0.97}
\definecolor{codecomment}{rgb}{0.35,0.45,0.35}
\definecolor{codekeyword}{rgb}{0.20,0.20,0.60}
\definecolor{codestring}{rgb}{0.55,0.25,0.25}

\lstset{
  basicstyle=\ttfamily\footnotesize,
  backgroundcolor=\color{codebg},
  commentstyle=\color{codecomment}\itshape,
  keywordstyle=\color{codekeyword}\bfseries,
  stringstyle=\color{codestring},
  frame=single,
  rulecolor=\color{black!30},
  breaklines=true,
  breakatwhitespace=true,
  columns=fullflexible,
  keepspaces=true,
  showstringspaces=false,
  tabsize=2,
  captionpos=b,
  xleftmargin=1em,
  xrightmargin=1em,
  aboveskip=1em,
  belowskip=1em,
}

% زبان‌های به‌کاررفته در این نوشتار
\lstdefinelanguage{promql}{
  keywords={sum,rate,count,increase,avg,max,min,by,without,offset},
  sensitive=true,
  comment=[l]{\#},
  morestring=[b]",
}
\lstdefinelanguage{yamlx}{
  keywords={true,false,null},
  sensitive=false,
  comment=[l]{\#},
  morestring=[b]",
}

% محیط آماده: قطعه‌کد چپ‌به‌راست
\lstnewenvironment{latincode}[1][]
  {\setLTR\lstset{#1}}
  {}

%-------------------------------------------------------------------------------
% ماکروهای کمکی این نوشتار
%-------------------------------------------------------------------------------
% یک شناسه‌ی لاتین در میان متن فارسی (نام بسته، نام تابع، نام بازو، …)
\DeclareRobustCommand{\code}[1]{\lr{\texttt{\small #1}}}
% یک عدد یا واحد لاتین در میان متن فارسی
\DeclareRobustCommand{\num}[1]{\lr{#1}}
% نام یک بازوی مقایسه
\DeclareRobustCommand{\arm}[1]{\lr{\texttt{\small #1}}}
% ارجاع به دفترچه‌ی آزمایشگاه (فقط در نسخه‌های پیش‌نویس؛ در متن نهایی استفاده نشده)
\DeclareRobustCommand{\notebookref}[1]{\textsuperscript{[د.آ. #1]}}

%-------------------------------------------------------------------------------
% سرصفحه‌های فهرست‌ها
%-------------------------------------------------------------------------------
\setlength{\cftbeforetoctitleskip}{-1.2em}
\setlength{\cftbeforelottitleskip}{-1.2em}
\setlength{\cftbeforeloftitleskip}{-1.2em}
\setlength{\cftaftertoctitleskip}{-1em}
\setlength{\cftafterlottitleskip}{-1em}
\setlength{\cftafterloftitleskip}{-1em}

\newcommand\tocheading{\par عنوان\hfill صفحه \par}
\newcommand\lofheading{\hspace*{.5cm}\figurename\hfill صفحه \par}
\newcommand\lotheading{\hspace*{.5cm}\tablename\hfill صفحه \par}

\renewcommand{\cftchapleader}{\cftdotfill{\cftdotsep}}
\renewcommand{\cfttoctitlefont}{\hspace*{\fill}\LARGE\bfseries}
\renewcommand{\cftaftertoctitle}{\hspace*{\fill}}
\renewcommand{\cftlottitlefont}{\hspace*{\fill}\LARGE\bfseries}
\renewcommand{\cftafterlottitle}{\hspace*{\fill}}
\renewcommand{\cftloftitlefont}{\hspace*{\fill}\LARGE\bfseries}
\renewcommand{\cftafterloftitle}{\hspace*{\fill}}

%-------------------------------------------------------------------------------
% قضیه‌ها، تعریف‌ها و نمونه‌ها
%-------------------------------------------------------------------------------
\theoremstyle{definition}
\newtheorem{definition}{تعریف}[section]
\newtheorem{remark}[definition]{نکته}
\newtheorem{note}[definition]{یادداشت}
\newtheorem{example}[definition]{نمونه}
\newtheorem{question}[definition]{سوال}
\newtheorem{remember}[definition]{یاداوری}
\theoremstyle{theorem}
\newtheorem{theorem}[definition]{قضیه}
\newtheorem{lemma}[definition]{لم}
\newtheorem{proposition}[definition]{گزاره}
\newtheorem{corollary}[definition]{نتیجه}

%-------------------------------------------------------------------------------
% صفحه‌های آغاز فصل
%-------------------------------------------------------------------------------
\makeatletter
\newcommand\mycustomraggedright{%
 \if@RTL\raggedleft%
 \else\raggedright%
 \fi}
\def\@makechapterhead#1{%
\thispagestyle{style1}
\vspace*{20\p@}%
{\parindent \z@ \mycustomraggedright
\ifnum \c@secnumdepth >\m@ne
\if@mainmatter

\bfseries{\Huge \@chapapp}\small\space {\chapternumber\thechapter}
\par\nobreak
\vskip 0\p@
\fi
\fi
\interlinepenalty\@M
\Huge \bfseries #1\par\nobreak
\vskip 120\p@

}
\newpage}
\bidi@patchcmd{\@makechapterhead}{\thechapter}{\tartibi{chapter}}{}{}
\bidi@patchcmd{\chaptermark}{\thechapter}{\tartibi{chapter}}{}{}
\makeatother

\pagestyle{fancy}
\renewcommand{\chaptermark}[1]{\markboth{\chaptername~\tartibi{chapter}: #1}{}}

\fancypagestyle{style1}{
\fancyhf{}
\fancyfoot[c]{\thepage}
\fancyhead[R]{\leftmark}%
\renewcommand{\headrulewidth}{1.2pt}
}

\fancypagestyle{style2}{
\fancyhf{}
\fancyhead[R]{چکیده}
\fancyfoot[C]{\thepage{}}
\renewcommand{\headrulewidth}{1.2pt}
}

\fancypagestyle{style3}{%
  \fancyhf{}%
  \fancyhead[R]{فهرست نمادها}
  \fancyfoot[C]{\thepage}%
  \renewcommand{\headrulewidth}{1.2pt}%
}

\fancypagestyle{style4}{%
  \fancyhf{}%
  \fancyhead[R]{فهرست جداول}
  \fancyfoot[C]{\thepage}%
  \renewcommand{\headrulewidth}{1.2pt}%
}

\fancypagestyle{style5}{%
  \fancyhf{}%
  \fancyhead[R]{فهرست اشکال}
  \fancyfoot[C]{\thepage}%
  \renewcommand{\headrulewidth}{1.2pt}%
}

\fancypagestyle{style6}{%
  \fancyhf{}%
  \fancyhead[R]{فهرست مطالب}
  \fancyfoot[C]{\thepage}%
  \renewcommand{\headrulewidth}{1.2pt}%
}

\fancypagestyle{style7}{%
  \fancyhf{}%
  \fancyhead[R]{نمایه}
  \fancyfoot[C]{\thepage}%
  \renewcommand{\headrulewidth}{1.2pt}%
}

\fancypagestyle{style8}{%
  \fancyhf{}%
  \fancyhead[R]{منابع و مراجع}
  \fancyfoot[C]{\thepage}%
  \renewcommand{\headrulewidth}{1.2pt}%
}
\fancypagestyle{style9}{%
  \fancyhf{}%
  \fancyhead[R]{واژه‌نامه‌ی فارسی به انگلیسی}
  \fancyfoot[C]{\thepage}%
  \renewcommand{\headrulewidth}{1.2pt}%
}

% حذف نام فهرست تصاویر و فهرست جداول از فهرست مطالب
\newcommand*{\BeginNoToc}{%
  \addtocontents{toc}{%
    \edef\protect\SavedTocDepth{\protect\the\protect\value{tocdepth}}%
  }%
  \addtocontents{toc}{%
    \protect\setcounter{tocdepth}{-10}%
  }%
}
\newcommand*{\EndNoToc}{%
  \addtocontents{toc}{%
    \protect\setcounter{tocdepth}{\protect\SavedTocDepth}%
  }%
}
\newcounter{savepage}
\renewcommand{\listfigurename}{فهرست اشکال}
\renewcommand{\listtablename}{فهرست جداول}
